\begin{center}
	\zihao{4}\heiti
	% 中国科学院大学\\
	\LARGE{$\boldsymbol{2007}$年硕士学位研究生入学统一考试试题\\
	电动力学A卷}

\end{center}
\noindent
考生须知：\\
1.本试卷满分为150分,全部考试时间总计180分钟。\\
2.所有答案必须写在答题纸上,写在试题纸上或草稿纸上一律无效。

\vspace*{2ex}
\hrule
\vspace*{4ex}

\fancypagestyle{mypagestyle}{%
	\fancyhf{}% Clear header/footer
	\fancyfoot[OR]{第~\thepage~页，共~\pageref{LastPage}~页}%
	\fancyfoot[EL]{第~\thepage~页，共~\pageref{LastPage}~页}%
	\fancyfoot[OL]{科目名称：电动力学 }%
	\fancyfoot[ER]{科目名称：电动力学 }%
	\renewcommand{\headrulewidth}{0pt}% Header rule of .4pt
	\renewcommand{\footrulewidth}{0.5pt}% Header rule of .4pt
}
\fancypagestyle{plain}{%
	\fancyhf{}% Clear header/footer
	\fancyfoot[OR]{第~\thepage~页，共~\pageref{LastPage}~页}%
	\fancyfoot[EL]{第~\thepage~页，共~\pageref{LastPage}~页}%
	\fancyfoot[OL]{科目名称：电动力学 }%
	\fancyfoot[ER]{科目名称：电动力学 }%
}

\setcounter{page}{1}%设置页面计数器
\pagestyle{mypagestyle}

\onehalfspacing
% 排版结束，下面是题目
\vspace{0.8cm}

\noindent
一、（共24分）问答题

\noindent
1）\hspace{0.2em}写出在介质中有源麦克斯韦方程组的积分形式，及解麦氏方程所必要的相关场量或源量之间的关系。
\vspace{0.2cm}

\noindent
2）\hspace{0.2em}平面电磁波从介电系数为$\large{\boldsymbol{\varepsilon_1}}$的透明介质以入射角$\boldsymbol{\theta}$（入射线与介质分界面法线的锐交角）入射到真空$\large{\boldsymbol{\varepsilon_2}}$ = $\large{\boldsymbol{\varepsilon_0}}$中（$\large{\boldsymbol{\varepsilon_1 > \varepsilon_0}}$），其E分量垂直于入射面。请写出电场强度和磁场强度的边界条件方程；假设发生了全反射，请写出全反射条件。
\vspace{0.2cm}

\noindent
3）\hspace{0.2em}写出电场和磁场的洛伦兹变换公式。

\vspace{0.5cm}
\noindent
二、（共36分）选择题

\noindent
1）\hspace{0.2em}在带自由面电荷的磁介质界面上，两边介质常数不同，这时候边值关系为：

\vspace{0.1cm}
a.\hspace{0.8em}磁感应强度法向不连续，磁场强度切向连续；

b.\hspace{0.8em}磁感应强度切向连续，磁场强度法向不连续；

c.\hspace{0.8em}磁感应强度法向连续，磁场强度切向不连续；

d.\hspace{0.8em}磁感应强度切向不连续，磁场强度法向连续。
\vspace{0.1cm}

\noindent
2）\hspace{0.2em}为探测波导管内$\mathrm{TE_{10}}$波的状态，需要在波导管壁某处开口以安插探测器，那么最佳开口位置应开在

\vspace{0.1cm}
a.\hspace{0.8em}窄边中部的纵向\hspace{3.2em}b.\hspace{0.2em}宽边中部的纵向

c.\hspace{0.8em}宽边中部的横向\hspace{3.2em}d.\hspace{0.2em}窄边中部的横向
\vspace{0.1cm}

% \begin{wrapfigure}{r}{-15pt}
% \includegraphics[width=3.6cm]{test/2007/pic_1.png} 
% \end{wrapfigure}

\noindent
3）\hspace{0.2em}如图所示，内半径为a，外半径为b的空心带电介质球壳，介电常数为$\large{\boldsymbol{\varepsilon}}$，球壳上带有均匀电荷，体密度为$\large{\boldsymbol{\rho}}$，则在空心内部A点（距球心的距离为
$\mathrm{r_1}$）和球壳内B点（距球心距离为
$\mathrm{r_2}$）的电场强度大小分别为

\vspace{0.2cm}
a.\hspace{0.8em}
$\large
{E_A = 0, \quad E_B = 
    \displaystyle
    {
        \frac{ \boldsymbol{\rho} } { 4\pi \boldsymbol{\varepsilon} r_2^2}
        \times \displaystyle
        { 
        \frac{4 \pi} {3} r_2^3
        } 
    }
}$

\vspace{0.2cm}
b.\hspace{0.8em}
$\large
{E_A = 
    \displaystyle
    {
        \frac{ \boldsymbol{\rho} } { 4\pi \boldsymbol{\varepsilon} r_1^2}
        \times \displaystyle
        { 
        (\frac{4 \pi} {3}  b^3 - \frac{4 \pi }{3} a^3 )
        } 
    }
    , \quad E_B = 
    \displaystyle
    {
        \frac{ \boldsymbol{\rho} } { 4\pi \boldsymbol{\varepsilon} r_2^2}
        \times \displaystyle
        { 
        \frac{4 \pi} {3} r_2^3
        } 
    }
}$

\vspace{0.3cm}
C.\hspace{0.8em}
$\large
{E_A = 0, \quad E_B = 
    \displaystyle
    {
        \frac{ \boldsymbol{\rho} } { 4\pi \boldsymbol{\varepsilon} r_2^2}
        \times \displaystyle
        { 
        (\frac{4 \pi} {3}  r_2^3 - \frac{4 \pi }{3} b^3)
        } 
    }
}$

\vspace{0.2cm}
d.\hspace{0.8em}
$\large
{E_A = 0, \quad E_B = 
    \displaystyle
    {
        \frac{ \boldsymbol{\rho} } { 4\pi \boldsymbol{\varepsilon} r_2^2}
        \times \displaystyle
        { 
        (\frac{4 \pi} {3} b^3 - \frac{4 \pi }{3} r_2^3)
        } 
    }
}$
\vspace{-4.8cm}
\begin{figure}[h!]
    \flushright
    \includegraphics[height=4.5cm]{test/2007/pic_1.png}
\end{figure}


\noindent
4）\hspace{0.2em}下面有关静电问题格林函数的说法中，哪一个是不正确的？

\vspace{0.1cm}
a.\hspace{0.8em}格林函数是单位点电荷产生的电位。

b.\hspace{0.8em}只要知道了点电荷的格林函数，任意连续分布电荷的电位原则上都可以求得。

c.\hspace{0.8em}在某些情况下，可以用电像法来求格林函数

d.\hspace{0.8em}格林函数满足拉普拉斯方程

\newpage
\noindent
5）\hspace{0.2em}
一个半径为a的电介质球，介电常数为$\large \boldsymbol{\varepsilon}$，极化强度为$\large \Vec{P} = p\sin \varphi \Vec{e_{\varphi}}$，其中$\varphi$为方位角。该球放置于$\Vec{E}=0$的环境中。则电介质球内的束缚电荷体密度$\rho_p$，球面上束缚电荷面密度$\sigma_p$在球坐标系$(r, \theta, \varphi)$中分别为

\vspace{0.2cm}
a.\hspace{0.8em}$\large \displaystyle {\rho_p = - \frac{p}{r \sin \theta} \cos \varphi, \hspace{0.4em} \sigma_p = 0 }$ \hspace{7em}
b.\hspace{0.8em}$\large \displaystyle {\rho_p = - \frac{p}{r \sin \theta} \cos \varphi, \hspace{0.4em} \sigma_p = p \sin \varphi }$

\vspace{0.4cm}
c.\hspace{0.8em}$\large \displaystyle {\rho_p = p \cos \varphi, \hspace{0.4em} \sigma_p = p \sin \varphi }$ \hspace{8em}
d.\hspace{0.8em}$\large \displaystyle {\rho_p = 0, \hspace{0.4em} \sigma_p = 0 }$

\vspace{0.4cm}
\noindent
6）\hspace{0.2em}$\nabla \cdot (\Vec{E} \times \Vec{H})$和$\Vec{H} \cdot \Vec{B}$在国际单位制中的单位分别为

\vspace{0.1cm}
a.\hspace{0.8em}瓦特/米$^2$和焦耳/米$^3$ \hspace{5em}b.\hspace{0.8em}瓦特/米$^3$和焦耳/米$^2$

\vspace{0.1cm}
a.\hspace{0.8em}焦耳/米$^3$和瓦特/米$^3$ \hspace{5em}b.\hspace{0.8em}瓦特/米$^3$和焦耳/米$^3$

\vspace{0.5cm}
\noindent
三、\hspace{0.2em}（30分）电荷Q放置在立体直角坐标$(0,\sqrt{2}a, \sqrt{2}a)$位置上，xz平面和xy平面都是无限大导电平板，另有一半径为a的导电圆球，圆心在原点上。求电四极距，$r \gg a$处的电势及电场的近似分布。这里$r = \displaystyle{ \sqrt{x^2 + y^2 + z^2} }$是观测点到原点的距离。
\vspace{-0.5cm}

\begin{figure*}[h]
    \flushright
        \subfloat[截面图]
        {\includegraphics[width=5cm]{test/2007/pic_2.png}}
        \subfloat[立体图]
        {\includegraphics[width=5cm]{test/2007/pic_3.png}}
\end{figure*}

\noindent
四、\hspace{0.2em}（30分）电荷量为-$e$质量为$m$的经典电子，在静电力的吸引作用下，初始以半径$a_0$绕质子作圆周旋转运动。如果没有量子效应，此系统将有电磁波辐射，消耗掉系统的能量。在辐射过程中，可把电子的运动过程看着是半径在逐渐缓慢减小的近似圆周旋转运动。

\noindent
1）\hspace{0.2em}这是什么辐射过程？

\noindent
2）\hspace{0.2em}写出辐射场（保留矢量积形式不做运算）和能流密度

\noindent
3）\hspace{0.2em}多长时间电子将落在核上？（不考虑各种更精细或更复杂的效应）

\vspace{3cm}
\noindent
五、\hspace{0.2em}（30分）静止长度为$l_1$的列车1和静止长度为$l_2$的列车2都呈现水平匀速直线运动状态$(l_1 > l_2)$，两列车平行运动且其后端在初始时刻重合对齐。现列车1以速度u相对于静止观测者前进，列车2又以速度v相对于列车1同向前进。求，在静止观测者看来，多长时间列车2的前端追上列车1的前端？

